% =============================================================================
% Chapter 17 -- Unified Conclusion
% =============================================================================

\chapter{Unified Conclusion}
\label{ch:conclusion}

This thesis has pursued a single coherent research programme across two interconnected works.
Parts~II--III developed a circular oversampling pipeline with BC-guided seed selection and three algorithms (GVM-CO, LRE-CO, LS-CO) that outperform linear interpolation methods in controlled benchmarks.
Part~IV tested the causal assumption underlying all oversampling research, using uniform random minority removal to isolate the effect of imbalance from confounding factors.
This chapter synthesises the findings, answers all research questions, identifies cross-cutting themes, and charts directions for future work.


% ============================================================================
\section{Summary of Contributions}
\label{sec:conc:contributions}

The thesis makes five primary contributions, spanning algorithmic design and empirical methodology.

\begin{enumerate}[label=\textbf{C\arabic*.}, leftmargin=*, itemsep=6pt]

  \item \textbf{BC-guided seed selection for minority-class oversampling.}
  Chapter~\ref{ch:seed_selection} introduces a four-stage pipeline (K-means clustering, PCA projection, BC + AGTP + JSD + $Z$ composite scoring, ranked selection) that identifies representative, geometrically faithful, and smooth seed sets.
  The composite criterion improves AGTP from 0.635 (random) to 0.839 and downstream F1 by $+0.025$ per classifier over ten benchmarks.

  \item \textbf{GVM-CO: gravity-biased Von Mises circular oversampling.}
  Chapter~\ref{ch:proposed_methods} introduces Gravity-biased Von Mises Circular Oversampling (Algorithm~\ref{alg:gvm_co}).
  The gravity score $g_c = \text{Density}_c / (\text{Spread}_c \cdot \text{Sparsity}_c + \varepsilon)$ weights each circle by its structural importance; the Von Mises direction $\mu = \arctan2(\mathbf{g}_c - \mathbf{c}_{ij})$ and concentration $\kappa = \kappa_{\max} \cdot s_{\text{combined}}^\gamma$ focus synthesis in the safest and densest regions.

  \item \textbf{LRE-CO and LS-CO: Voronoi-constrained and layered-segmental oversampling.}
  LRE-CO (Local Region Estimation, Algorithm~\ref{alg:lre_co}) partitions the minority class into Voronoi cells and enforces a certainty threshold ($\tau = 0.80$) that rejects circles crossing cluster boundaries.
  LS-CO (Layered Segmental, Algorithm~\ref{alg:ls_co}) decomposes each circle into $L$ concentric rings and $M$ angular segments, directing synthesis toward the cluster's gravity center via exponential layer weights.
  Together these three algorithms cover qualitatively different geometric regimes of minority class structure.

  \item \textbf{A controlled degradation-and-recovery protocol for causal imbalance analysis.}
  Chapter~\ref{ch:causality_intro} introduces a two-phase experimental protocol.
  Phase~1 removes minority instances one step at a time using uniform random selection, ensuring no distributional bias so that performance drops are causally attributable to the imbalance ratio itself.
  Phase~2 applies seven oversamplers incrementally from the most-degraded state to full balance.
  The four-rule LeakageSafeCV protocol prevents the three forms of data leakage identified by Kaufman et al.~\cite{kaufman2012leakage}.

  \item \textbf{Empirical causal evidence: imbalance degrades performance; recovery is incomplete.}
  Chapter~\ref{ch:degradation} confirms that performance degrades monotonically across all five classifiers and six metrics, with a non-linear threshold at 40--50\% removal ($\IR \approx 2$).
  Chapter~\ref{ch:recovery} shows that no oversampler achieves full recovery: the best methods (GVM-CO, KM-SMOTE) reach 98.4\% of the balanced AUC, leaving a residual gap of 1.4\% attributable to irreducible information loss.

\end{enumerate}


% ============================================================================
\section{Answers to Research Questions}
\label{sec:conc:rqs}

\begin{table}[ht]
\centering
\caption{Answers to the seven research questions (RQ1--RQ7).}
\label{tab:conc:rqs}
\small
\begin{tabular}{@{}p{0.7cm}p{12.5cm}@{}}
\toprule
\textbf{RQ} & \textbf{Answer} \\
\midrule
RQ1 & Composite $\Phi$ criterion raises mean AGTP by $+0.204$ vs.\ random seeds; downstream F1 $+0.025$, Sensitivity $+0.035$. PCA projection is essential (AGTP drops $-0.211$ without). \emph{(Ch.~\ref{ch:seed_selection})} \\[3pt]
RQ2 & Disk-based generation yields F1 $+0.002$--$+0.013$ over SMOTE; largest gains at $\IR>9$ (LRE-CO $+0.018$). \emph{(Ch.~\ref{ch:proposed_methods}--\ref{ch:results})} \\[3pt]
RQ3 & Friedman significant for G-Mean ($p=0.0002$) and BAcc ($p=0.0001$); not for F1 ($p=0.6251$) or AUC ($p=0.1139$). Study underpowered ($\approx30\%$) at $N=14$; need $N\geq45$ for 80\% power. SVM-SMOTE leads on G-Mean. \emph{(Ch.~\ref{ch:statistical})} \\[3pt]
RQ4 & Seed selection is the dominant component ($+0.053$ F1; AGTP most important). K-means $>$ HAC. Optimal: $K=3$--$4$, $\kappa_{\max}=20$, $L=10$, $\tau=0.80$. Native-dimension mode $+0.002$--$+0.003$ F1 at high $d$. \emph{(Ch.~\ref{ch:ablation})} \\[3pt]
RQ5 & \textbf{Yes.} Monotonic degradation with shape preserved; non-linear threshold at 40--50\% removal ($\IR\approx2$). \emph{(Ch.~\ref{ch:degradation})} \\[3pt]
RQ6 & RF most robust ($\Delta$AUC $=0.191$); KNN/LR most sensitive. Hierarchy: RF$>$MLP$>$DT$>$LR$\approx$KNN. \emph{(Ch.~\ref{ch:degradation})} \\[3pt]
RQ7 & No full recovery. GVM-CO $\approx$ KM-SMOTE ceiling: AUC 0.871 ($\Delta=-0.014$). Residual gap is irreducible when the surviving minority is sparse. \emph{(Ch.~\ref{ch:recovery}--\ref{ch:discussion})} \\
\bottomrule
\end{tabular}
\end{table}


% ============================================================================
\section{Cross-Cutting Findings}
\label{sec:conc:crosscutting}

Three themes run through all works and emerge most clearly in their synthesis.

\paragraph{1. LRE-CO is the most consistently robust circular oversampler.}
Across 14 benchmark datasets, five classifiers, six metrics, and an independent causal recovery study, LRE-CO achieves the best or joint-best performance in the largest number of settings.
Its Voronoi constraint (Equation~\eqref{eq:voronoi}) and certainty threshold ($\tau = 0.80$) prevent the boundary leakage that limits standard SMOTE and B-SMOTE.
GVM-CO shows selective advantages at intermediate imbalance ratios and high-dimensional datasets; LS-CO provides stable mid-tier performance with low computational cost.

\paragraph{2. Imbalance is a genuine cause of degradation, but the gap is irreducible.}
The causal study (Part~IV) confirms that imbalance \emph{per se} degrades classifiers, not merely the confounds of distributional shift or reduced sample size.
However, the recovery experiment reveals a fundamental limit: synthetic oversampling recovers 90--97\% of the balanced baseline, with a residual gap attributable to the inherent information loss when fewer real instances are available.
This finding has a direct practical implication: at high imbalance ratios, the priority should shift from oversampling strategy to data collection, since no algorithm can overcome irreducible information loss.

\paragraph{3. Leakage is quantifiable and consequential.}
The LeakageSafeCV protocol prevents 2--5\% AUC inflation that would otherwise be attributable to oversampling leakage rather than genuine improvement.
This effect is large enough to reverse the relative ranking of oversampling methods and to produce the false impression of statistical significance in underpowered studies.
The protocol should be treated as a prerequisite for reproducible benchmarking in the oversampling literature.


% ============================================================================
\section{Limitations}
\label{sec:conc:limitations}

\paragraph{Circular oversampling (Parts~II--III).}
The PCA projection discards information above two dimensions; native-dimension mode avoids this but requires Von Mises--Fisher approximation whose accuracy degrades in very high dimensions ($d > 100$).
Hyperparameter tuning ($K$, $\kappa_{\max}$, $L$, $\tau$) adds complexity; the optimistic-bias disclosure (Chapter~\ref{ch:experimental_setup}) acknowledges that reported numbers reflect tuned configurations.
The benchmark is limited to 14 datasets; statistical power is insufficient ($\approx 30\%$) for rank-2 effect detection.

\paragraph{Causality study (Part~IV).}
Random removal is inherently stochastic; different random sequences may produce slightly different degradation curves. We mitigate this with five-fold cross-validation and averaging, but exact removal order effects are not characterized.
Synthetic dataset analogues replace unavailable KEEL/UCI files; qualitative conclusions are expected to generalise but exact numbers may differ.
The study covers only $\IR \leq 4$; behavior at extreme imbalance ($\IR > 10$) is not characterized.
Only binary classification is studied.


% ============================================================================
\section{Future Work}
\label{sec:conc:future}

\paragraph{Circular oversampling.}
\begin{itemize}[leftmargin=*]
  \item Extend evaluation to $N \geq 45$ datasets to achieve 80\% statistical power for rank-2 effect detection.
  \item Investigate the native-dimension Von Mises--Fisher approximation at $d > 100$ and develop a more accurate high-dimensional direction sampler.
  \item Apply LRE-CO and GVM-CO to multi-class imbalance using one-vs-rest decomposition.
  \item Combine the circular framework with deep generative models (CTGAN, TVAE) to handle tabular data with complex dependencies.
\end{itemize}

\paragraph{Causal study extensions.}
\begin{itemize}[leftmargin=*]
  \item Extend the degradation sweep to $\IR$ up to 50 to characterize the full performance curve at extreme imbalance~\cite{he2009learning}.
  \item Validate on real KEEL/UCI datasets with matching $(n, d)$ to confirm the synthetic-analogue findings~\cite{fernandez2018learning}.
  \item Investigate multi-class imbalance, where the interaction between ratios across classes adds complexity not addressed here~\cite{haixiang2017learning}.
  \item Design an oversampler that explicitly targets the irreducible information-loss gap using uncertainty quantification, prioritising synthesis in the most uncertain minority regions~\cite{he2008adasyn}.
\end{itemize}

\paragraph{Unified future direction.}
The most promising synthesis is an adaptive oversampling system that uses distributional monitoring (via the Bhattacharyya Coefficient) to continuously assess minority-class coverage in production data streams, triggering targeted data collection or adaptive oversampling when coverage drops below a threshold.
GVM-CO provides the synthesis engine; the causal study quantifies the performance recovery achievable at each coverage level.
Together, they form a closed-loop framework for managing class imbalance in live systems.
